\documentclass[12pt]{article}
\usepackage[T1]{fontenc}
\usepackage[utf8]{inputenc}
\usepackage{lmodern}
\usepackage{textcomp}
\usepackage{enumitem}
\usepackage{geometry}
\usepackage{graphicx}
\usepackage{amsmath, amssymb}
\geometry{margin=1in}
\begin{document}
\begin{center}
Psychology - Undergraduate Level Assessment 16 \
Total Marks: 40 \
Answer all questions.
\end{center}

\section*{Multiple Choice (10 marks)}
\begin{enumerate}[label=\arabic*.]
\item The reticular formation and reticular activating system are associated with all of the following functions, EXCEPT:
\begin{enumerate}[label=(\alph*)]
    \item Alertness
    \item Sensation of taste
    \item Hunger and thirst regulation
    \item Decussation of auditory stimuli
    \item Pain
\end{enumerate}
\item The diagnosis of Alzheimer's disease can be most accurately verified by
\begin{enumerate}[label=(\alph*)]
    \item Blood test showing elevated levels of certain proteins
    \item Electroencephalogram (EEG) showing abnormal brain activity
    \item ruling out other etiologies through a comprehensive psychodiagnostics workup
    \item neuropsychological testing showing the expected profile of deficits and spared functions
    \item Physical examination and patient history
\end{enumerate}
\item When subjected to moderate punishment, an instrumental response of moderate strength will
\begin{enumerate}[label=(\alph*)]
    \item be temporarily suppressed but strengthen over time
    \item be eliminated
    \item be strengthened
    \item remain unchanged
    \item increase in frequency
\end{enumerate}
\item The more difference shown by the behavior of identical twins raised apart, the more the differences in their behavior can be attributed to their
\begin{enumerate}[label=(\alph*)]
    \item hormonal changes
    \item innate abilities
    \item genetic traits
    \item environments
    \item educational backgrounds
\end{enumerate}
\item Mr. Gordon suffered damage to the back of his right frontal lobe. As a result, he is unable to
\begin{enumerate}[label=(\alph*)]
    \item Recognize faces
    \item move his left hand
    \item Hear sounds correctly
    \item understand information he reads
    \item speak intelligibly
\end{enumerate}
\end{enumerate}

\section*{True / False (10 marks)}
\textit{Answer True or False}
\begin{enumerate}[label=\arabic*.]
\setcounter{enumi}{5}
\item True or False: Learned helplessness is an example of cognitive dissonance, where individuals feel powerless to change their circumstances.
\item True or False: The best presentation time to produce classical conditioning is a continuous pairing of the conditioned and unconditioned stimuli.
\item True or False: Object relations therapy primarily focuses on the client's conscious thoughts and feelings rather than early childhood experiences.
\item True or False: A 66-year-old client with similar symptoms is more likely to be diagnosed with Alzheimer's disease than Parkinson's disease.
\item True or False: It is advisable to ignore ethical violations among psychologists, as they are personal matters that don't require intervention.
\end{enumerate}

\section*{Long Form Response (20 marks)}
\textit{Answer each question with a clear and complete explanation.}
\begin{enumerate}[label=\arabic*.]
\setcounter{enumi}{10}
\item Discuss the evidence supporting Incremental Theory and One-trial Theory in learning, specifically focusing on their applicability in cognitive versus motor tasks.
\item Discuss the significance of providing education and job training to adolescents and young adults post-drug treatment in terms of secondary intervention strategies in psychology.
\end{enumerate}

\vfill
\noindent\textit{End of Paper}
\end{document}